%----------------------------------------------------------------------------------------
% HYPOTHESETOETSEN. Alle opgaven zijn eigen werk.
%
% Opbouw: Start, leerstof (theorie, werkwijze of uitleg, voorbeeld, korte oefening), en op een
% nieuwe pagina de Opdrachten (door elkaar, stijgende moeilijkheid), Koppel en Test jezelf.
% Labels van dit hoofdstuk beginnen met h3-.
%----------------------------------------------------------------------------------------

\mychapter{3}{Hypothesetoetsen}

\lkr{Leerplan B+S'': LPD 54 (hypothesen toetsen met nulhypothese, alternatieve hypothese, p-waarde en significantieniveau), met de wenken over de steekproevenverdeling, eenzijdige en tweezijdige toetsen, toetsen voor proporties en gemiddelden, de twee soorten fouten en de centrale limietstelling. Toetsen voor een proportie doen we exact met de binomiale verdeling; toetsen voor een gemiddelde met een bekende $\sigma$ en de normale verdeling. De t-toets valt buiten het leerplan. De p-waarde mag met ICT.}

Een producent beweert dat hoogstens 2~\% van zijn producten defect is. Je test er 200 en vindt er 9 defect. Liegt hij, of had je pech? Met een hypothesetoets beslis je dat op een vaste manier. Je gaat ervan uit dat de bewering klopt, en berekent hoe waarschijnlijk je resultaat dan is. Is dat heel onwaarschijnlijk, dan verwerp je de bewering. In dit hoofdstuk leer je die redenering, die in alle wetenschappen gebruikt wordt, en de valkuilen ervan.

\begin{oefeningen}[Start]
\oef{h3-start-1}

\oef{h3-verken}
\end{oefeningen}

\lkr{Laat oefening \ref{oef:h3-verken} maken vóór de uitleg. De redenering van vraag c (\textit{als de dobbelsteen eerlijk is, hoe groot is dan de kans op zoveel zessen?}) is de kern van elke toets.}

%========================================================================================
\section{De steekproevenverdeling}

In hoofdstuk 1 zag je dat een steekproefgrootheid schommelt van steekproef tot steekproef. Om te beoordelen of een resultaat verrassend is, moet je weten hoe sterk.

\begin{theorie}
De \textbf{\textcolor{accent}{steekproevenverdeling}} van een steekproefgrootheid is haar kansverdeling over alle mogelijke aselecte steekproeven van grootte $n$.

\begin{stelling}{Steekproevenverdeling van een proportie en een gemiddelde}
Voor aselecte steekproeven van grootte $n$ uit een grote populatie is

\vspace{1mm}
\qquad $\formule{\E(\hat{p}) = p}$, \quad $\formule{\sigma(\hat{p}) = \sqrt{\dfrac{p(1 - p)}{n}}}$ \qquad en \qquad $\formule{\E(\overline{X}) = \mu}$, \quad $\formule{\sigma(\overline{X}) = \dfrac{\sigma}{\sqrt{n}}}$

\vspace{1mm}
\textbf{Centrale limietstelling}: voor grote $n$ is $\overline{X}$ bij benadering normaal verdeeld, welke verdeling de populatie ook heeft.
\end{stelling}
\end{theorie}

De formule voor $\hat{p}$ volgt uit de binomiale verdeling: $\hat{p} = \frac{X}{n}$ met $X \sim B(n, p)$ (zie de Koppel van hoofdstuk 1). De formules voor $\overline{X}$ en de centrale limietstelling bewijzen we niet. Vier keer zoveel waarnemingen halveert de spreiding van het gemiddelde: daarom geven grote steekproeven nauwkeurigere schattingen.

\naast{\begin{minipage}{\linewidth}
De figuur toont de steekproevenverdeling van $\overline{X}$ voor een populatie met $\mu = 70$ en $\sigma = 8$, voor steekproeven van 4, 16 en 64. De standaardafwijking is $\frac{8}{\sqrt n}$: 4, 2 en 1. Hoe groter de steekproef, hoe dichter de gemiddelden bij 70 liggen.
\end{minipage}}{\begin{tikzpicture}[x=.36cm,y=7cm,font=\scriptsize]
  \draw[->,grijs] (56,0) -- (84.5,0);
  \foreach \t in {60,65,70,75,80}{\node[below] at (\t,0) {\t};}
  \draw[grijs,thick,dashed] plot[domain=57:83,samples=80,smooth] (\x,{gauss(\x,70,4)});
  \draw[accent,thick] plot[domain=60:80,samples=80,smooth] (\x,{gauss(\x,70,2)});
  \draw[fel!80!black,thick] plot[domain=64:76,samples=80,smooth] (\x,{gauss(\x,70,1)});
  \node[fel!80!black,right] at (71.2,.38) {$n = 64$};
  \node[text=accent,right] at (72.5,.17) {$n = 16$};
  \node[right] at (75.5,.06) {$n = 4$};
\end{tikzpicture}}

\voorbeeld{}{In een stad doet 30~\% aan vrijwilligerswerk. Wat is de kans dat in een aselecte steekproef van 100 inwoners minstens 40~\% vrijwilligerswerk doet?\label{vb:h3-steekproef}}{%
\uitleg{We gebruiken de steekproevenverdeling van $\hat{p}$, die voor $n = 100$ bij benadering normaal is.}
\begin{stappen}
\sigma(\hat{p}) = \sqrt{\frac{0{,}3 \cdot 0{,}7}{100}} \approx 0{,}046 & Met $\E(\hat{p}) = 0{,}3$. \\[3ex]
P(\hat{p} \geqslant 0{,}4) \approx 0{,}015 & \texttt{normalcdf(0.4,1E99,0.3,0.0458)}: ongeveer 1,5~\%.
\end{stappen}
Een steekproef met 40~\% vrijwilligers is dus zeldzaam als $p = 0{,}3$. Vind je ze toch, dan twijfel je aan $p = 0{,}3$. Dat is het idee van een toets.}

\begin{oefeningen}
\oef{h3-steekproef}
\end{oefeningen}

%========================================================================================
\section{De toets}

\subsection{Begrippen}

\begin{theorie}
De \textbf{\textcolor{accent}{nulhypothese}} $\formule{H_0}$ is de bewering die je toetst, met een vaste waarde van de parameter, zoals $p = \frac16$. De \textbf{\textcolor{accent}{alternatieve hypothese}} $\formule{H_1}$ is wat je vermoedt als $H_0$ niet klopt, zoals $p > \frac16$.

\vspace{1mm}
De \textbf{\textcolor{accent}{p-waarde}} is de kans, als $H_0$ waar is, op een steekproefresultaat dat minstens zo extreem is als het waargenomen resultaat (in de richting van $H_1$).

\vspace{1mm}
Het \textbf{\textcolor{accent}{significantieniveau}} $\formule{\alpha}$ kies je vooraf, meestal $0{,}05$. Is de p-waarde $\leqslant \alpha$, dan \textbf{\textcolor{accent}{verwerp}} je $H_0$: het resultaat is \textbf{\textcolor{accent}{significant}}. Anders verwerp je $H_0$ niet.
\end{theorie}

Een toets is een bewijs uit het ongerijmde met kansen. Je neemt aan dat $H_0$ klopt. Leidt dat tot een resultaat dat bijna nooit gebeurt, dan verwerp je die aanname. Belangrijk: als je $H_0$ niet verwerpt, heb je niet bewezen dat ze waar is. Je hebt alleen te weinig aanwijzingen tegen haar gevonden.

\subsection{Een toets voor een proportie}

\voorbeeld{}{Een speler vermoedt dat een dobbelsteen te vaak zes gooit. Bij 120 worpen valt 30 keer zes. Toets bij $\alpha = 0{,}05$.\label{vb:h3-proportie}}{%
\uitleg{Bij een eerlijke dobbelsteen is $p = \frac16$, en verwacht je 20 zessen. We berekenen hoe waarschijnlijk 30 of meer zessen zijn als de dobbelsteen eerlijk is.}
\begin{stappenbreed}
H_0\colon p = \tfrac16 \qquad H_1\colon p > \tfrac16 & De speler vermoedt te veel zessen: een eenzijdige toets. \\[1ex]
\text{onder } H_0\colon \ X \sim B(120,\, \tfrac16) & $X$ is het aantal zessen. \\[1ex]
\text{p-waarde} = P(X \geqslant 30) = 1 - P(X \leqslant 29) \approx 0{,}013 & Met \texttt{binomcdf(120,1/6,29)}. \\[1ex]
0{,}013 \leqslant 0{,}05: \ H_0 \text{ verwerpen} & Significant: de dobbelsteen gooit waarschijnlijk te vaak zes.
\end{stappenbreed}}

\begin{werkwijze}
\textbf{1.}~Formuleer $H_0$ (met $=$) en $H_1$ ($>$, $<$ of $\neq$), en kies $\alpha$. \ \textbf{2.}~Bepaal de verdeling van de toetsingsgrootheid als $H_0$ waar is. \ \textbf{3.}~Bereken de p-waarde: de kans op het waargenomen resultaat of iets extremer, in de richting van $H_1$. \ \textbf{4.}~Vergelijk met $\alpha$. \ \textbf{5.}~Formuleer het besluit in de context, voorzichtig.
\end{werkwijze}

\subsection{Een toets voor een gemiddelde}

Bij een gemiddelde gebruik je de steekproevenverdeling van $\overline{X}$. In dit hoofdstuk kennen we de standaardafwijking $\sigma$ van de populatie, bijvoorbeeld uit vroegere metingen.

\voorbeeld{}{Een machine vult flessen met gemiddeld 500~ml, met $\sigma = 2$~ml. Een consumentenorganisatie meet 25 flessen en vindt een gemiddelde van 499,1~ml. Vult de machine gemiddeld te weinig? Toets bij $\alpha = 0{,}05$.\label{vb:h3-gemiddelde}}{%
\uitleg{Als de machine goed werkt, is $\overline{X}$ normaal verdeeld met gemiddelde 500 en standaardafwijking $\frac{2}{\sqrt{25}} = 0{,}4$. We berekenen de kans op een gemiddelde van 499,1 of minder.}
\begin{stappenbreed}
H_0\colon \mu = 500 \qquad H_1\colon \mu < 500 & Het vermoeden: te weinig. \\[1ex]
\text{onder } H_0\colon \ \overline{X} \sim N(500;\, 0{,}4) & $\sigma(\overline{X}) = \frac{2}{\sqrt{25}}$. \\[1ex]
\text{p-waarde} = P(\overline{X} \leqslant 499{,}1) \approx 0{,}012 & $z = \frac{499{,}1 - 500}{0{,}4} = -2{,}25$. \\[1ex]
0{,}012 \leqslant 0{,}05: \ H_0 \text{ verwerpen} & De machine vult waarschijnlijk gemiddeld te weinig.
\end{stappenbreed}}

\nujij{Bij een toets is de p-waarde 0,08. Wat besluit je bij $\alpha = 0{,}05$? Betekent dit dat $H_0$ waar is?}{$0{,}08 > 0{,}05$: $H_0$ niet verwerpen. Dat bewijst niet dat $H_0$ waar is; de aanwijzingen ertegen zijn alleen niet sterk genoeg.}

\begin{oefeningen}
\oef{h3-toets}
\end{oefeningen}

%========================================================================================
\section{Eenzijdig en tweezijdig}

Vermoed je alleen dat een parameter \textit{anders} is, zonder richting, dan is $H_1$ van de vorm $p \neq p_0$. Dan telt een extreem resultaat aan beide kanten. De p-waarde is dan het dubbel van de kans in de staart aan de kant van het resultaat.

\voorbeeld{}{Een munt landt 60 keer op kop bij 100 worpen. Is de munt oneerlijk? Toets tweezijdig bij $\alpha = 0{,}05$.\label{vb:h3-tweezijdig}}{%
\uitleg{\textit{Oneerlijk} kan naar kop of naar munt: $H_1\colon p \neq 0{,}5$. Het resultaat ligt aan de kant van te veel kop; we verdubbelen die staart.}
\begin{stappenbreed}
H_0\colon p = 0{,}5 \qquad H_1\colon p \neq 0{,}5 & Tweezijdig. \\[1ex]
P(X \geqslant 60) \approx 0{,}028 \quad \Rightarrow \quad \text{p-waarde} \approx 2 \cdot 0{,}028 = 0{,}057 & Met $X \sim B(100;\, 0{,}5)$; 40 of minder keer kop is even extreem. \\[1ex]
0{,}057 > 0{,}05: \ H_0 \text{ niet verwerpen} & Niet significant.
\end{stappenbreed}}

Merk op: had je vooraf vermoed dat de munt te vaak kop geeft, dan was de p-waarde 0,028 en had je $H_0$ verworpen. Daarom kies je $H_1$ \textit{voor} je de data bekijkt, op basis van de vraag, en niet op basis van het resultaat.

\begin{oefeningen}
\oef{h3-tweezijdig}
\end{oefeningen}

%========================================================================================
\section{Twee soorten fouten}

Een toets beslist op basis van een steekproef, en kan zich dus vergissen.

\begin{theorie}
{\centering\arrayrulecolor{accent}\renewcommand{\arraystretch}{1.6}%
\begin{tabular}{|>{\centering\arraybackslash}m{3.6cm}|>{\centering\arraybackslash}m{4.4cm}|>{\centering\arraybackslash}m{4.4cm}|}
\hline
 & \textbf{$H_0$ is waar} & \textbf{$H_0$ is niet waar} \\ \hline
\textbf{$H_0$ verwerpen} & \textbf{\textcolor{accent}{fout van de eerste soort}} (type I) & juist \\ \hline
\textbf{$H_0$ niet verwerpen} & juist & \textbf{\textcolor{accent}{fout van de tweede soort}} (type II) \\ \hline
\end{tabular}\par}

\vspace{1mm}
De kans op een fout van de eerste soort is hoogstens $\formule{\alpha}$.
\end{theorie}

Kies je $\alpha$ kleiner, dan verwerp je minder snel ten onrechte, maar mis je vaker een echt effect: de kans op een fout van de tweede soort stijgt. Een grotere steekproef verkleint die tweede fout. Bij een medische test is een fout van de eerste soort bijvoorbeeld: een werkzaam lijkend medicijn dat niet werkt; van de tweede soort: een werkzaam medicijn dat je afwijst.

\voorbeeld{}{Bij 20 worpen met een munt verwerp je $H_0\colon p = 0{,}5$ ten voordele van $H_1\colon p > 0{,}5$ als er minstens 15 keer kop valt. Wat is de kans op een fout van de eerste soort? En op een fout van de tweede soort als in werkelijkheid $p = 0{,}7$?\label{vb:h3-fouten}}{%
\uitleg{Een fout van de eerste soort: $H_0$ is waar, en toch valt 15 of meer keer kop. Een fout van de tweede soort: $p = 0{,}7$, en toch valt minder dan 15 keer kop.}
\begin{stappenbreed}
P(X \geqslant 15 \mid p = 0{,}5) \approx 0{,}021 & De kans op een fout van de eerste soort: kleiner dan 0,05. \\[1ex]
P(X \leqslant 14 \mid p = 0{,}7) \approx 0{,}58 & Bij $p = 0{,}7$ mis je de afwijking in 58~\% van de gevallen.
\end{stappenbreed}
Met maar 20 worpen ontdek je een munt die 70~\% kop geeft dus maar in 42~\% van de gevallen.}

\begin{oefeningen}
\oef{h3-fouten}
\end{oefeningen}

%========================================================================================
\opdrachtenkop

\begin{oefeningen}
\oef{h3-o1}

\oef{h3-o2}

\oef{h3-o3}

\oef{h3-o4}

\oef{h3-o5}

\oef{h3-o6}

\oef{h3-o7}

\oef{h3-o8}

\oef{h3-o9}
\end{oefeningen}

%----------------------------------------------------------------------------------------
% KOPPEL: verbanden met andere leerstof
%----------------------------------------------------------------------------------------
\needspace{20\baselineskip}
\opdrachtsectie{Koppel}

\lkr{De eerste oefening vergelijkt een toets met een bewijs uit het ongerijmde (Logica en bewijzen, 6de jaar). De tweede gebruikt voorwaardelijke kansen uit Kansrekenen.}

\begin{oefeningen}
\oef{h3-koppel-1}

\oef{h3-koppel-2}
\end{oefeningen}

%----------------------------------------------------------------------------------------
\opdrachtsectie{Test jezelf}

\begin{oefeningen}
\oef{h3-t1}

\oef{h3-t2}

\oef{h3-t3}

\oef{h3-t4}
\end{oefeningen}

\herhaalschema{%
  \fv{oef:h3-t1}{vb:h3-proportie}{oef:h3-toets}%
  \fv{oef:h3-t2}{vb:h3-gemiddelde}{oef:h3-toets}%
  \fv{oef:h3-t3}{vb:h3-tweezijdig}{oef:h3-tweezijdig}%
  \fzv{oef:h3-t4}{oef:h3-fouten}}

\ifzelfstudie
\section*{Maak je eigen samenvatting}
Sluit dit boek en schrijf in je schrift: de steekproevenverdeling van een proportie en van een gemiddelde, de centrale limietstelling, de begrippen nulhypothese, alternatieve hypothese, p-waarde en significantieniveau, de vijf stappen van een toets, het verschil tussen eenzijdig en tweezijdig, en de twee soorten fouten. Schrijf ook op wat het betekent als je $H_0$ niet verwerpt. Vergelijk daarna met de kaders.
\fi
